% Generated draft: requires mathematical and source review.
\documentclass[11pt]{article}
\usepackage[margin=25mm]{geometry}
\usepackage{fontspec}
\setmainfont{Latin Modern Roman}
\usepackage{amsmath,amssymb,mathtools,unicode-math}
\setmathfont{Latin Modern Math}
\usepackage{enumitem,needspace,xurl,hyperref}
\setlength{\parindent}{0pt}
\setlength{\parskip}{5pt}
\setlength{\emergencystretch}{3em}
\newcommand{\N}{\mathbb{N}}
\newcommand{\Z}{\mathbb{Z}}
\newcommand{\Q}{\mathbb{Q}}
\newcommand{\R}{\mathbb{R}}
\newcommand{\C}{\mathbb{C}}
\newcommand{\F}{\mathbb{F}}
\newcommand{\A}{\mathbb{A}}
\newcommand{\PP}{\mathbb{P}}
\newcommand{\E}{\mathbb{E}}
\newcommand{\eps}{\varepsilon}
\newcommand{\dd}{\,\mathrm{d}}
\begin{document}
\section*{UnsolvedMath ID: 30002713}
\textbf{Source problem code:} catalogue entry
\subsection{Definitions and mathematical statement}
Let \(p\) be an odd prime and let \(n\geq 1\) be an integer. Define the dihedral group
\[
D_{p^n}=\langle r,s\mid r^{p^n}=s^2=1,\ srs^{-1}=r^{-1}\rangle,
\qquad |D_{p^n}|=2p^n.
\]
The restriction \(n\geq 1\) is the conventional positive prime-power interpretation; the supplied statement does not explicitly specify the range of \(n\).

For a precise global interpretation, call a finite group \(G\) an \emph{Oort group in characteristic \(p\)} if the following lifting property holds. For every algebraically closed field \(k\) of characteristic \(p\), every smooth projective connected curve \(X/k\), and every faithful action \(\rho:G\hookrightarrow\operatorname{Aut}_k(X)\), there exist a complete discrete valuation ring \(R\) of characteristic zero with residue field identified with \(k\), a smooth proper relative curve \(\mathcal X/R\), and an action of \(G\) by \(R\)-automorphisms of \(\mathcal X\), together with a \(G\)-equivariant isomorphism \(\mathcal X\times_R k\cong X\) recovering \(\rho\).

\emph{Question under this global interpretation.} Is \(D_{p^n}\) an Oort group in characteristic \(p\) for every odd prime \(p\) and every integer \(n\geq 1\)?

There is an unresolved terminology issue in the supplied record: its question says ``Oort groups,'' whereas its literature assessment says ``local Oort groups.'' For comparison, a finite group \(G\) is a \emph{local Oort group in characteristic \(p\)} if, for every algebraically closed field \(k\) of characteristic \(p\) and every faithful continuous action \(G\hookrightarrow\operatorname{Aut}_k(k[[t]])\), there exist a complete discrete valuation ring \(R\) of characteristic zero with residue field \(k\) and a continuous action of \(G\) by \(R\)-algebra automorphisms of \(R[[T]]\) whose reduction is equivariantly isomorphic to the given action on \(k[[t]]\). The corresponding local question asks whether this property holds for \(D_{p^n}\) for every odd prime \(p\) and every \(n\geq 1\). These two formulations are recorded separately; the supplied evidence does not identify which definition was intended or establish an equivalence between them.
\subsection{Short English statement}
The question asks whether every dihedral group of order \(2p^n\), for an odd prime \(p\) and a positive integer \(n\), has the Oort lifting property from characteristic \(p\) to characteristic zero. The source question uses global terminology, while the supplied literature assessment discusses local actions on a formal disc; the intended formulation requires review.
\subsection{Sources}
\begin{itemize}
\item Supplied local problem registry record, permanent UnsolvedMath ID \(30002713\), source problem code \texttt{catalogue entry}, titled \emph{Odd Dihedral Groups as Oort Groups}; collection: Oberwolfach Reports Open Problems. The supplied proposed year is 2014; no individual proposer is identified.
\item Supplied report citation: \emph{Valuation Theory and Its Applications} (2015), Oberwolfach Reports, pp.~2757--2823, DOI 10.4171/owr/2014/49. \url{https://doi.org/10.4171/owr/2014/49}. The background describes Oberwolfach Reports as workshop reports documenting mathematical discussions and open problems.
\item Supplied literature source: \emph{Oberwolfach Report 49/2014}. \url{https://ems.press/content/serial-article-files/46540}.
\item Supplied literature source: Obus, \emph{A generalization of the Oort conjecture}. \url{https://arxiv.org/abs/1502.07623}.
\item Literature assessment dated exactly \texttt{2026-08-21}: the supplied curation labels the problem partially solved and states that the local Oort conjecture for all odd-prime-power dihedral groups has substantial special cases, including \(D_p\) and \(D_9\), but remains open in general. This is supplied dated curation derived from the OWR record, not an independently verified assessment. The source explicitly requests independent verification of time-sensitive status and cited claims.
\item The registry labels the statement \texttt{corrected\_verified} and reports a manual semantic, typographic, and LaTeX review of its cleanup. This supplied editorial label does not establish mathematical correctness or resolve the global/local terminology issue.
\end{itemize}
\end{document}
